% !TEX TS-program = pdflatex
\documentclass[10pt,a4paper]{altacv}

\geometry{left=1.2cm,right=1.2cm,top=0.8cm,bottom=0.6cm,footskip=0pt}

\usepackage[T1]{fontenc}
\usepackage[utf8]{inputenc}
\usepackage[french]{babel}
\usepackage[default]{lato}
\usepackage{microtype}
\DisableLigatures{encoding=*, family=*} % Désactive les ligatures pour compatibilité ATS maximale

\definecolor{Navy}{HTML}{0A2540}
\definecolor{SlateGrey}{HTML}{2E2E2E}
\definecolor{LightGrey}{HTML}{555555}
\colorlet{heading}{Navy}
\colorlet{accent}{Navy}
\colorlet{emphasis}{SlateGrey}
\colorlet{body}{LightGrey}

\renewcommand{\itemmarker}{{\small\textbullet}}
\renewcommand{\ratingmarker}{\faCircle}

\usepackage[hidelinks]{hyperref}
\hypersetup{
  pdftitle={CV -- Achraf Izelmadene -- Candidature ETF VINCI Construction},
  pdfauthor={Achraf Izelmadene},
  pdfsubject={Stage PFE Ingénieur Fiabilisation Équipement SIMEF},
  pdfkeywords={ETF, VINCI, SIMEF, freinage, fiabilisation, maintenance industrielle, TIA Portal, AMDEC, stage PFE}
}

% ── Style sections compactes ─────────────────────────────────────────
\renewcommand{\cvsection}[2][]{%
  \vspace{3pt}%
  {\color{heading}\Large\bfseries #2 \hspace{0.2em}\leaders\vrule height 3pt depth -1.5pt\hfill\mbox{}}\par\vspace{2.5pt}
}

\renewcommand{\cvevent}[4]{%
  {\normalsize\bfseries\color{accent}#1\par}\vspace{1pt}
  \ifstrequal{#2}{}{}{%
    {\small\itshape\color{emphasis}#2\ifstrequal{#4}{}{}{ - #4}}%
  }%
  ~\hfill
  \ifstrequal{#3}{}{}{{\small\color{emphasis}#3}}%
  \par\vspace{2pt}
}

\newenvironment{cvitems}{\begin{itemize}[leftmargin=*, label=\textbullet, nosep, itemsep=1.5pt]}{\end{itemize}}
\newcommand{\cvprojet}[1]{{\normalsize\bfseries\color{accent} #1}\par\vspace{1.5pt}}

% ── En-tête ATS centré + rubrique À PROPOS ───────────────────────────
\makeatletter
\newcommand{\target}[1]{\def\@target{#1}}
\newcommand{\apropos}[1]{\def\@apropos{#1}}

\renewcommand{\makecvheader}{%
  \begingroup
    \centering\color{emphasis}%
    {\LARGE\bfseries\@name\par}\vspace{2.5pt}
    {\normalsize\bfseries\color{accent}\@tagline\par}\vspace{3.5pt}
    {\footnotesize\color{emphasis}\@personalinfo\par}
  \endgroup\par\vspace{2pt}
  \ifboolexpr{test {\ifdef{\@target}} or test {\ifdef{\@apropos}}}{%
    \cvsection{À PROPOS}
    \ifdef{\@target}{{\small\textbf{\@target}\par}\vspace{2pt}}{}%
    \ifdef{\@apropos}{{\small\color{emphasis}\@apropos\par}}{}%
  }{}%
  \vspace{1pt}
}
\makeatother

\begin{document}

\name{Izelmadene Achraf}
\tagline{Élève Ingénieur en Génie Électrique et Systèmes Intelligents}
\personalinfo{%
  \href{mailto:achraf.izelmadene@gmail.com}{achraf.izelmadene@gmail.com}%
  \hspace{0.9em}\textcolor{body!60}{|}\hspace{0.9em}%
  +212699098736%
  \hspace{0.9em}\textcolor{body!60}{|}\hspace{0.9em}%
  Marrakech, Maroc (Mobilité géographique)%
  \hspace{0.9em}\textcolor{body!60}{|}\hspace{0.9em}%
  Permis B%
}

\target{Stage de fin d’études (PFE) de 6 mois dès février 2027 : Ingénieur Fiabilisation Équipement chez ETF (AMM).}
\apropos{Élève ingénieur en dernière année, orienté fiabilité des systèmes industriels, automatisme et contrôle. Expérience terrain en unité de production (maintenance préventive/corrective, migration automate TIA Portal) et méthodologie d’industrialisation (cahier des charges, analyse des risques AMDEC, instrumentation, suivi budgétaire).}

\makecvheader

\cvsection{FORMATION}
\cvevent{Étudiant en cycle d’ingénieur en Génie Électrique et Systèmes Intelligents}{École Marocaine des Sciences de l’Ingénieur - Marrakech}{2025 - Présent}{}
{\footnotesize\color{emphasis}Production \& Maintenance (préventive \& prédictive), Lean Management, Management \& Contrôle qualité, Hygiène, Sécurité \& Environnement (HSE), Supervision industrielle \& Digital Twin, Smart capteurs \& actionneurs, Régulation industrielle, Gestion de projet...}\par\vspace{2.5pt}
\cvevent{Licence SMP, option Automatique}{Faculté des Sciences Semlalia - Marrakech}{2020 - 2024}{}
{\footnotesize\color{emphasis}Asservissements linéaires, Automatismes et capteurs, Traitement du signal, Électronique analogique...}\par

\cvsection{EXPÉRIENCE}
\cvevent{Stagiaire en Unité de Production}{Menara Prefa}{juil. 2024}{Marrakech}
\begin{cvitems}
  \item Maintenance préventive et corrective sur lignes automatisées, suivi des indicateurs d’équipements et de disponibilité.
  \item Participation à la migration d’un automate vers Siemens S7 (TIA Portal) et mise à jour de la documentation technique.
\end{cvitems}
\vspace{2pt}
\cvevent{Stagiaire Bureau d’Études}{Core Engineering Solutions}{juil. - août 2026}{Berrchid, Casablanca}
\begin{cvitems}
  \item Dimensionnement complet d’un poste de livraison HTA/BT, bilan de puissance et conformité aux normes techniques.
  \item Étude d’intégration PV en autoconsommation, schémas unifilaires et notes de calcul (AutoCAD, Caneco BT, PVsyst).
\end{cvitems}

\cvsection{PROJETS}
\cvprojet{Contrôle Qualité Intelligent \& Banc d’Inspection (Vision \& Électronique Industrielle)}
\begin{cvitems}
  \item Étude complète d’un système d’inspection automatisée : cahier des charges, analyse des risques (AMDEC) et instrumentation.
  \item Dimensionnement financier du projet (calcul de VAN, ROI, budget) et rédaction des modes opératoires de contrôle.
\end{cvitems}
\vspace{2pt}
\cvprojet{Maintenance Prédictive 4.0 par AutoML (TPOT vs. Auto-WEKA)}
\begin{cvitems}
  \item Détection anticipée de défaillances sur données capteurs (vibrations, température, pression) et benchmark F1-score/Recall.
  \item Définition de seuils critiques d’alerte complétée par une étude CND (Contrôle Non Destructif) par analyse d’huile.
\end{cvitems}
\vspace{2pt}
\cvprojet{Détection d’anomalies Smartgrid (Clustering \& Supervision)}
\begin{cvitems}
  \item Benchmark d’algorithmes de clustering (KMeans, DBSCAN, GMM) sur des flux de données IoT de réseaux électriques.
  \item Dashboard interactif de visualisation des clusters pour le monitoring d’anomalies et l’optimisation des contrôles.
\end{cvitems}

\cvsection{COMPÉTENCES}
{\small
\begin{tabular}{@{}ll@{}}
\textbf{Maintenance \& Fiabilité :} & Préventive \& corrective, Indicateurs (MTBF/MTTR), AMDEC, CND, Modes opératoires \\[1.5pt]
\textbf{Automatisme \& Équipements :} & TIA Portal, Bancs de test, Ladder, Grafcet, SCADA / HMI, Modbus RTU, Caneco BT \\[1.5pt]
\textbf{Outils, Budget \& Data :} & Suivi budgétaire (VAN/ROI), Excel (suivi \& calculs), Power BI, Python (Pandas), Pack Office \\[1.5pt]
\textbf{Qualités personnelles :} & Travail d’équipe, Rigueur \& ponctualité, Esprit d’entreprendre, Curiosité technique
\end{tabular}
}

\cvsection{CERTIFICATIONS}
\begin{cvitems}
  \item \textbf{Foundations of IoT Systems and Industrial Automation}~\hfill {\color{emphasis}LearnQuest (Coursera)}
  \item \textbf{Sensors and Sensor Circuit Design}~\hfill {\color{emphasis}University of Colorado Boulder (Coursera)}
  \item \textbf{Certified 3DEXPERIENCE Collaborative Industry Innovator}~\hfill {\color{emphasis}Dassault Systèmes}
  \item \textbf{MATLAB / Simulink Onramp}~\hfill {\color{emphasis}MathWorks}
\end{cvitems}

\cvsection{ACTIVITÉS \& ENGAGEMENT}
\begin{cvitems}
  \item \textbf{Hack The Intelligence (AISEC 2026) - 1\textsuperscript{er} Prix :} Développement d’un prototype Smart Agriculture IoT.
  \item \textbf{Vision AI Competition (EMSI 2026) - 2\textsuperscript{ème} Prix :} Système de traduction en langue des signes (ASL).
  \item \textbf{Congrès COCIA’26 - Présentation Scientifique :} Prototype d’éclairage public piézoélectrique SmartStruct.
  \item \textbf{Engagement Associatif \& Bénévolat (2025) :} Soutien scolaire et encadrement pédagogique d’enfants orphelins.
\end{cvitems}

\cvsection{LANGUES}
\textbf{Arabe :} Langue maternelle~\hfill \textbf{Français :} DELF B2~\hfill \textbf{Anglais :} Professionnel

\end{document}
